\section{An exact obstruction-space criterion for a pure Weil character}
\label{sec:pure-weil-complexes}
This section turns strong semiregularity for a pure Weil character into an exact dimension test. It then constructs such a character on a Gaussian eightfold and proves why a natural family of complexes representing it fails the test.

Let $A$ have dimension $g=2n\ge4$, let $K$ act with signature $(n,n)$, and write $W_K\subset H^{2n}(A,\Q)$ for its Weil plane. For a perfect complex $E$, its Hochschild semiregularity map is
\[
 \sigma_E^2:\operatorname{Ext}^2(E,E)\longrightarrow HH_{-2}(A).
\]
If a finite group $G\subset A\times A^\vee$ acts coherently by translations and tensor products with degree-zero line bundles, and $E$ is $G$-equivariant, call $E$ \emph{strongly $G$-semiregular} when this map is injective on the invariant self-Ext space. The coherence and the equivariant structure are part of the hypothesis. A finite group of other automorphisms is not covered by this definition here.

\begin{theorem}[The exact dimension test]\label{thm:pure-weil-ext}
Suppose that $B\in H^{1,1}(A)\cap H^2(A,\Q)$ and
\[
 e^B\operatorname{ch}(E)=w\in W_K\setminus\{0\}.
\]
Then, for the preceding action and equivariant structure,
\[
 E\text{ is strongly }G\text{-semiregular}
 \quad\Longleftrightarrow\quad
 \dim_{\C}\operatorname{Ext}^2(E,E)^G
     =2\binom{2n}{2}=2n(2n-1).
\]
The assertion includes the trivial group. Every such equivariant object has invariant second self-Ext dimension at least $2n(2n-1)$. Neither simplicity nor vanishing of negative self-Ext is needed for this equivalence. The latter vanishing is a separate hypothesis for the deformation theorem.
\end{theorem}

\begin{proof}
Put $T=H^0(A,T_A)$, $U_0=H^1(A,\mathcal O_A)$ and $H=U_0\oplus T$. The tangent bundle is trivial and the Todd class is one. Multiplicative Hochschild--Kostant--Rosenberg identifies
\[
 HH^*(A)=\bigwedge^*H,\qquad \dim H=2g,
\]
where $(\alpha,\xi)$ acts on Hodge cohomology by $\alpha\wedge-+\iota_\xi$. The characteristic maps
\[
 \operatorname{ev}_E^j:HH^j(A)\longrightarrow\operatorname{Ext}^j(E,E)
\]
form a graded algebra homomorphism, and the trace identity gives
\begin{equation}\label{eq:pure-weil-trace}
 \sigma_E^j\operatorname{ev}_E^j(x)=x\cdot\operatorname{ch}(E).
\end{equation}
We use the full Hochschild action, including the bivector and gerbe summands; the relevant Atiyah--Chern identities are \cite[Proposition~6.2.1, Corollary~6.3.2 and Theorem~6.4.2]{BFHH}.

Choose a nonzero holomorphic volume form. Cap product with it identifies $HH^j(A)$ with $HH_{g-j}(A)$, so the character corresponds to an element of $\bigwedge^gH$. Serre duality pairs $\operatorname{Ext}^2(E,E)$ perfectly with $\operatorname{Ext}^{g-2}(E,E)$. Under this pairing and the Hochschild trace pairing, $\sigma_E^2$ is adjoint to $\operatorname{ev}_E^{g-2}$. Indeed pairing the categorical trace of an endomorphism $u$ with a Hochschild operation $x$ is the Serre trace of $u\operatorname{ev}_E(x)$. Equivalently, the same identity follows term by term from the expression $\operatorname{tr}(u\operatorname{At}(E)^j/j!)$.

Translations and degree-zero tensor products act trivially on Hochschild cohomology and preserve the volume and Serre trace. Evaluation therefore lands in invariant self-Ext. Averaging a dual vector shows that Serre duality restricts to a perfect pairing on invariants. If $R^j=\operatorname{Ext}^j(E,E)^G$, the adjunction becomes
\begin{equation}\label{eq:pure-weil-dual}
 \sigma_E^2|_{R^2}\text{ injective}
 \quad\Longleftrightarrow\quad
 HH^{g-2}(A)\longrightarrow R^{g-2}\text{ surjective}.
\end{equation}

The rational $B$-field is handled without choosing a line bundle representing it. Define the invertible shear
\[
 T_B(\alpha,\xi)=(\alpha+\iota_\xi B,\xi).
\]
The type $(1,1)$ of $B$ places the first component in $U_0$. Directly,
\[
 x\cdot(e^Bz)=e^B(T_Bx\cdot z).
\]
Thus $\operatorname{ev}_{E,B}=\operatorname{ev}_E\circ\bigwedge T_B$ is still an algebra homomorphism, $\sigma_{E,B}=e^B\sigma_E$ has the same injectivity, and their composite is cap product with $w$. Surjectivity in \eqref{eq:pure-weil-dual} is unchanged by this shear.

Write $H^{1,0}(A)=X\oplus Y$ for the two $K$-eigenspaces, each of dimension $n$. Then
\[
 H^1_\sigma=X\oplus\overline Y,
 \qquad H^1_{\bar\sigma}=Y\oplus\overline X.
\]
The rational nonzero class $w$ has nonzero coefficients on both conjugate determinant lines. Under the volume identification it becomes
\[
 \eta=a\operatorname{vol}_{H_1}+b\operatorname{vol}_{H_2},
 \quad ab\ne0,\qquad
 H_1=Y^*\oplus\overline Y,\quad H_2=X^*\oplus\overline X.
\]
Here $H=H_1\oplus H_2$ and $\dim H_i=g$. Multiplication by $\eta$ on $\bigwedge^2H$ has kernel exactly $H_1\otimes H_2$ and rank $2\binom g2$. The mixed products kill both top forms. The pure second exterior powers multiply injectively by the opposite top form, with distinct bidegrees since $g\ge4$.

The composite in \eqref{eq:pure-weil-trace}, after the shear, consequently has rank $2\binom g2$. This proves the lower bound on $\dim R^2$ and proves sufficiency when equality holds.

For necessity suppose that the semiregularity map on $R^2$ is injective. Every element of its quadratic annihilator $H_1\otimes H_2$ evaluates to zero, by \eqref{eq:pure-weil-trace}. Evaluation is an algebra homomorphism, so it kills the exterior ideal generated by these mixed products. Its degree-$g-2$ image is therefore a quotient of
\[
 \bigwedge^{g-2}H_1\oplus\bigwedge^{g-2}H_2,
\]
of dimension $2\binom g2$. Equation~\eqref{eq:pure-weil-dual} says that this image is all of $R^{g-2}$. Invariant Serre duality gives $\dim R^2=\dim R^{g-2}\le2\binom g2$, and the reverse inequality was already proved. This proves necessity and the theorem.
\end{proof}

Strong semiregularity implies the weak equivariant condition in Perry's theorem: the forgetful semiregularity square commutes, so an element killed by the equivariant trace is also killed by the injective ordinary trace on invariants. No general equality of the weak and strong definitions is needed here. For a pure Weil character on an eightfold, the exact strong criterion is $\dim\operatorname{Ext}^2(E,E)^G=56$.

\subsection{An integral character on the principal Gaussian family}

Let $(X,h)$ be a principally polarized abelian fourfold and set $A=X\times X$. The endomorphism $\rho(x,y)=(-y,x)$ gives a Gaussian action of signature $(4,4)$ preserving the product polarization $\theta$. With respect to a symplectic basis of $H^1(X,\Z)$, the Gaussian Hermitian lattice has four blocks $\left(\begin{smallmatrix}0&i\\-i&0\end{smallmatrix}\right)$. Multiplication of a basis vector in each block by a unit gives the hyperbolic plane. Thus this point belongs to the principal family $U^4$.

Let $m$ be addition and put
\[
 S=p_1^*h-p_2^*h,\qquad
 \ell=m^*h-p_1^*h-p_2^*h,\qquad
 \Omega=(S+i\ell)/2.
\]
For $a=x+iy\in\Z[i]$ put $D(a)=xS-y\ell=a\Omega+\bar a\bar\Omega$, an integral divisor class. Choose its line bundle multiplicatively in $a$. If $u_1,\ldots,u_8$ is a symplectic basis on the first factor and $v_1,\ldots,v_8$ the corresponding basis on the second, then $\zeta_j=u_j+iv_j$ is an $i$-eigenvector and
\[
 \Omega=\frac12\sum_{j=1}^4\zeta_{2j-1}\zeta_{2j}.
\]
It is nondegenerate on the eight-dimensional eigenspace, with $\Omega^5=0$ and $\bar\Omega^5=0$.

The Hermitian matrix of $D(a)$, relative to the product polarization, is
\[
 \begin{pmatrix}x&-y\\-y&-x\end{pmatrix}\otimes I_4.
\]
For $a\ne0$ its eigenvalues are $|a|$ and $-|a|$, each four times. The index theorem for line bundles gives cohomology only in degree four, unchanged by a degree-zero twist, with dimension $\Nm(a)^4$. This instance can also be seen by an isogeny: rational congruence diagonalizes the indicated nonsingular rational symmetric matrix to one positive and one negative rational entry. After clearing denominators the corresponding isogeny pulls the line bundle back to an exterior product of an ample bundle on $X$ and an antiample bundle on $X$, up to a degree-zero twist. Kodaira vanishing, Serre duality and K\"unneth give concentration in degree four. Pullback by a finite \'etale isogeny contains the original cohomology as a direct summand, so the same concentration holds before pullback. Riemann--Roch gives the stated dimension.

\begin{proposition}[A pure sixteen-term character]\label{prop:sixteen-line-character}
Fix $M\ge2$ and set
\[
 (c_1,c_2,c_3,c_4)=(1,i,M(1+i),M(1-i)),
 \quad a_J=\sum_{j\in J}c_j,\quad L_J=L_{D(a_J)}.
\]
The sixteen classes $D(a_J)$ are distinct. For
\[
 \xi=\sum_{J\subset\{1,2,3,4\}}(-1)^{|J|}[L_J],
 \qquad B=-\frac12D(c_1+c_2+c_3+c_4),
\]
one has the exact identity
\[
 e^B\operatorname{ch}(\xi)
   =w=2iM^2(\Omega^4-\bar\Omega^4)\ne0.
\]
This is a rational Weil class, with no other Chern-character components.
\end{proposition}

\begin{proof}
Equality of two subset sums would give
\[
 \epsilon_1+M(\epsilon_3+\epsilon_4)=0,
 \qquad \epsilon_2+M(\epsilon_3-\epsilon_4)=0,
 \qquad \epsilon_j\in\{-1,0,1\}.
\]
Since $M\ge2$, both parenthesized integers and then all $\epsilon_j$ vanish. The divisor map is injective, proving distinctness.

The alternating exponential is
\[
 e^B\operatorname{ch}(\xi)
 =\prod_{j=1}^4 2\sinh(D(c_j)/2).
\]
Its arguments are $S,-\ell,M(S-\ell),M(S+\ell)$, so its degree-four component is
\[
 -M^2S\ell(S^2-\ell^2)=2iM^2(\Omega^4-\bar\Omega^4).
\]
Only codimensions four, six and eight can occur. The codimension-six term is $w\sum_jD(c_j)^2/24$. But
\[
 \sum_jD(c_j)^2=(1+2M^2)(S^2+\ell^2)
   =4(1+2M^2)\Omega\bar\Omega,
\]
whose product with $w$ vanishes. Finally $(x,y)\mapsto(x,-y)$ fixes $S$ and negates $\ell$. It negates the whole hyperbolic-sine product, while acting as $+1$ on top cohomology. Its codimension-eight term is therefore zero.

The two fourth powers span the conjugate Weil lines. They are independent, so $w\ne0$; its displayed expression as a product of integral divisors proves rationality. More precisely, if $w_1=\operatorname{Im}(\zeta_1\cdots\zeta_8)$, then $\Omega^4=(3/2)\zeta_1\cdots\zeta_8$ and $w=-6M^2w_1$. With the product principal orientation,
\[
 \int_A w_1^2=128,\qquad \int_Aw^2=4608M^4.
\]
Indeed the product of the two conjugate volume elements integrates to $256$, while their individual squares vanish. Also $\theta w=0$, since each summand of $\theta$ contains one factor from each Gaussian eigenspace.
\end{proof}

\subsection{Every two-level gluing has the same excessive obstruction space}

The next proposition rules out the whole extension family, including its simple members. Put
\[
 P=\bigoplus_{|J|\text{ even}}L_J,
 \qquad Q=\bigoplus_{|J|\text{ odd}}L_J.
\]

\begin{proposition}\label{prop:sixteen-line-gluing}
For every $\delta\in\operatorname{Ext}^4(P,Q)$, let $E_\delta$ be defined by the triangle
\[
 Q[3]\longrightarrow E_\delta\longrightarrow P
       \xrightarrow{\delta}Q[4].
\]
Then $E_\delta$ has no negative self-Ext and has normalized character $w$ from Proposition~\ref{prop:sixteen-line-character}. It satisfies
\[
 \operatorname{Ext}^2(E_\delta,E_\delta)
       \simeq H^2(A,\mathcal O_A)^{\oplus16},
 \qquad\dim\operatorname{Ext}^2(E_\delta,E_\delta)=448.
\]
Every finite translation and degree-zero-tensor action for which it is equivariant fixes this entire space. Consequently none of these complexes is strongly equivariantly semiregular. If all sixty-four components of $\delta$ are nonzero, the complex is simple nevertheless.
\end{proposition}

\begin{proof}
The triangle is a cone construction and has precisely two cohomology sheaves, $H^0(E_\delta)=P$ and $H^{-3}(E_\delta)=Q$. Its $K$-class is $[P]-[Q]=\xi$. Distinct line bundles in the construction have differences of index four, so
\[
 \operatorname{Ext}^k(L_J,L_{J'})=0\quad(J\ne J',\ k\ne4),
 \qquad\operatorname{Ext}^k(L_J,L_J)=H^k(A,\mathcal O_A).
\]
Use the self-hyper-Ext spectral sequence
\[
 E_2^{p,q}=\bigoplus_i
 \operatorname{Ext}^p(H^i(E_\delta),H^{i+q}(E_\delta))
 \ \Longrightarrow\ \operatorname{Ext}^{p+q}(E_\delta,E_\delta).
\]
Its only rows are $q=0,\pm3$. Each off-diagonal row is supported at $p=4$. Thus there are no terms of negative total degree, proving gluability. In total degree two the sole term is
\[
 E_2^{2,0}=\operatorname{Ext}^2(P,P)\oplus\operatorname{Ext}^2(Q,Q)
       =H^2(A,\mathcal O_A)^{\oplus16}.
\]
The only possible outgoing differential is $d_4$ to $\operatorname{Ext}^6(P,Q)=0$; the only possible incoming differential starts at $\operatorname{Ext}^{-2}(Q,P)=0$. There are no other rows that can reach or leave this term. The edge map is the asserted isomorphism. Its dimension is $16\binom82=448$.

In degree zero, the $d_4$ differential from diagonal scalar endomorphisms sends the scalar at one endpoint of a component of $\delta$ minus the scalar at its other endpoint to that multiple of the component. If all components are nonzero, the complete bipartite graph is connected and forces all sixteen scalars to agree. There is no other term in total degree zero, proving simplicity.

An equivariant translation--tensor functor preserves cohomology and first Chern classes. Distinctness of the sixteen classes therefore prevents it from permuting their line summands. On each $\operatorname{Ext}^2(L_J,L_J)=H^2(\mathcal O_A)$, translation acts trivially and the scalar ambiguity of the identification acts trivially by conjugation. The edge isomorphism is equivariant, so every one of the $448$ classes is invariant. Theorem~\ref{thm:pure-weil-ext} now excludes strong semiregularity.
\end{proof}

These actions also fail the weak equivariant condition. The projection of an allowed group fixing $E_\delta$ to the translation factor $A$ is injective: a tensor functor alone fixing one $L_J$ has trivial line bundle. Tensor conjugation by one fixed $L_J^{-1}$ makes the action a free finite translation action, possibly with a scalar coherence cocycle. Descent identifies its invariant category with a twisted derived category on the quotient abelian variety. On Hochschild homology the forgetful map is, under local Morita identifications and the invertible tensor conjugation, pullback along the finite \'etale isogeny. This is injective over $\C$. The semiregularity square therefore has the same kernel before and after forgetting equivariance, and the weak condition fails as well.

The calculation explains a necessary change in the construction. For a two-level triangle $Q[r]\to E\to P\to Q[r+1]$, the diagonal second self-Ext can only be affected by
\[
 \operatorname{Ext}^{1-r}(Q,P)\longrightarrow
 \operatorname{Ext}^2(P,P)\oplus\operatorname{Ext}^2(Q,Q)
 \longrightarrow\operatorname{Ext}^{r+3}(P,Q).
\]
For $r=3$ these two groups vanish in the example. Changing only the gluing maps cannot repair it. The shifts $3|J|$ introduce a further problem: a gap at least six gives a mixed $\operatorname{Ext}^4$ term of negative total degree. No differential can kill it, because a differential from a diagonal row would begin in negative Ext degree and all other mixed rows are supported only in Ext degree four. Thus that choice of more cohomological levels is not gluable.

Proposition~\ref{prop:sixteen-line-character} gives an explicit class on the principal hyperbolic family, but no semiregular representative of it. A positive application of Perry's theorem would require a different gluable perfect representative, with coherent allowed equivariance and invariant second self-Ext dimension $56$. The entire normalized character would then remain Hodge along the Gaussian Weil family. This is a concrete existence problem for a complex, beyond the verified character identity.
